\chapter*{Определения, обозначения и сокращения}
\addcontentsline{toc}{chapter}{Определения, обозначения и сокращения}

\begin{description}

\item[Автоэнкодер (Autoencoder)] — нейронная сеть, обучающаяся сжимать данные в латентное представление и восстанавливать их обратно.

\item[Абляционное исследование (Ablation study)] — эксперимент, в котором отдельные компоненты модели или входных данных отключаются или заменяются, чтобы оценить их вклад в результат.

\item[Авторегрессионная модель (Autoregressive model)] — модель, которая генерирует последовательность по шагам, предсказывая следующий токен на основе уже сгенерированного контекста.

\item[Аккорд (Chord)] — одновременное звучание трёх и более нот, образующих гармоническую структуру.

\item[Baseline] — базовый вариант сравнения, относительно которого оценивается качество разработанной модели.

\item[Биосонификация (Biosonification)] — преобразование биологических данных в звуковую форму для восприятия человеком.

\item[Bio conditioning] — использование биологического вектора или профиля признаков как управляющего условия для нейросетевой генерации.

\item[Bio→Harmony→Melody] — иерархическая схема генерации, в которой биологические признаки сначала управляют построением гармонии, а затем мелодия генерируется с учётом этой гармонии.

\item[Вариационный автоэнкодер (VAE, Variational Autoencoder)] — генеративная модель, использующая вероятностное латентное пространство.

\item[Генеративно-состязательная сеть (GAN, Generative Adversarial Network)] — архитектура из двух сетей (генератор и дискриминатор), обучающихся в состязательном режиме.

\item[Гармония (Harmony)] — аккордовая структура музыкальной композиции.

\item[Градиентное накопление (Gradient Accumulation)] — техника обучения, позволяющая эмулировать большой batch size на ограниченной памяти.

\item[Взаимная информация (Mutual Information)] — мера статистической зависимости между двумя случайными величинами; показывает, насколько знание одной величины уменьшает неопределённость другой.

\item[Checkpoint] — сохранённое состояние обученной модели, которое можно загрузить для продолжения обучения или генерации.

\item[Chord-tone ratio] — доля мелодических нот, принадлежащих текущему аккорду; используется как метрика согласованности мелодии и гармонии.

\item[Control profile] — компактный вектор управляющих признаков, описывающий целевые музыкальные свойства фрагмента, например темп, плотность нот и регистр.

\item[Cross-entropy loss] — функция потерь для задачи предсказания следующего токена; показывает, насколько вероятностное распределение модели отличается от правильного ответа.

\item[Длинная краткосрочная память (LSTM, Long Short-Term Memory)] — тип рекуррентной нейронной сети с механизмом управления памятью.

\item[Embedding] — численное векторное представление объекта, например биологического фрагмента или токена, пригодное для обработки нейросетью.

\item[Epoch] — один полный проход алгоритма обучения по обучающей выборке.

\item[Изоэлектрическая точка (pI, Isoelectric Point)] — pH, при котором белок имеет нулевой суммарный заряд.

\item[Кодон (Codon)] — тройка нуклеотидов, кодирующая одну аминокислоту.

\item[Loss] — численная функция ошибки, оптимизируемая при обучении модели.

\item[MAESTRO] — датасет фортепианных MIDI-записей, часто используемый в исследованиях генерации символической музыки.

\item[Мелодия (Melody)] — последовательность нот, образующая музыкальную линию.

\item[Механизм внимания (Attention Mechanism)] — компонент нейронной сети, позволяющий модели фокусироваться на релевантных частях входных данных.

\item[Memory tokens] — служебные токены, через которые биологический embedding передаётся трансформеру как дополнительный контекст.

\item[MIDI (Musical Instrument Digital Interface)] — стандарт цифрового представления музыки в виде событий (note on/off, pitch, velocity).

\item[Mixed Precision] — техника обучения нейронных сетей с использованием 16-битных и 32-битных чисел для экономии памяти.

\item[MusicLang] — высокоуровневый текстовый язык описания музыки через аккорды, тональности и относительные ноты.

\item[Нуклеотид (Nucleotide)] — структурная единица ДНК или РНК, состоящая из азотистого основания, сахара и фосфатной группы.

\item[Октава (Octave)] — интервал между двумя нотами, частота которых отличается в два раза.

\item[Открытая рамка считывания (ORF, Open Reading Frame)] — участок последовательности между стартовым и стоп-кодонами.

\item[Полутон (Semitone)] — минимальный интервал в западной музыкальной системе (1/12 октавы).

\item[Рекуррентная нейронная сеть (RNN, Recurrent Neural Network)] — нейронная сеть, обрабатывающая последовательности с использованием внутреннего состояния.

\item[Real bio] — вариант абляционного сравнения, в котором модель получает настоящий биологический embedding и настоящий управляющий профиль.

\item[RefSeq] — база эталонных биологических последовательностей, используемая как источник геномных данных.

\item[Seed] — фиксированное начальное значение генератора случайных чисел, обеспечивающее воспроизводимость эксперимента.

\item[Shuffled bio] — вариант абляционного сравнения, в котором биологическое условие перемешивается между фрагментами для проверки вклада связи ``биоданные — музыка''.

\item[Символическая музыка (Symbolic Music)] — представление музыки в виде дискретных событий (нот), а не аудиосигнала.

\item[Сонификация (Sonification)] — преобразование данных в звуковую форму для восприятия или анализа.

\item[Neutral bio] — вариант абляционного сравнения, в котором биологическое условие заменяется нейтральным или усреднённым представлением.

\item[Такт (Bar, Measure)] — единица музыкального времени, содержащая фиксированное количество долей.

\item[Токенизация (Tokenization)] — преобразование данных в последовательность дискретных токенов для обработки моделью.

\item[Top-k matching] — подбор одного из k ближайших музыкальных сегментов по расстоянию между биологическим управляющим профилем и музыкальными дескрипторами.

\item[Train/validation/test split] — разбиение данных на обучающую, валидационную и тестовую части для обучения, настройки и независимой проверки модели.

\item[Трансформер (Transformer)] — архитектура нейронной сети, основанная на механизме внимания, без рекуррентных связей.

\item[Validation loss] — значение функции потерь на валидационной выборке, которая не используется для обновления параметров модели.

\item[Valid MIDI rate] — доля сгенерированных MIDI-файлов, которые корректно открываются и содержат ожидаемую музыкальную структуру.

\item[Энтропия последовательности (Sequence Entropy)] — мера неопределённости или разнообразия символов в последовательности.

\item[FASTA] — текстовый формат для представления биологических последовательностей (DNA, RNA, protein).

\item[GC-content] — доля гуанина и цитозина в нуклеотидной последовательности.

\item[GRAVY (Grand Average of Hydropathy)] — средняя гидрофобность аминокислот в белковой последовательности.

\item[k-mer] — подпоследовательность длины k (например, триплет нуклеотидов при k=3).

\item[POP909] — датасет из 909 популярных песен в MIDI-формате с аннотациями аккордов и мелодий.

\item[PyTorch] — фреймворк глубокого обучения с динамическими вычислительными графами.

\end{description}
